\section{From Language Chaos to Verified Merge}

Software history is often presented as a clean ladder of progress: assembly to C, C to managed runtimes, dynamic languages to typed product stacks, and now AI-assisted programming. The lived history is messier. Programming languages are human inventions shaped by fashion, hardware pressure, academic taste, corporate platforms, tooling accidents, hiring markets, and the local pain of the decade. Each new language promises to compress some burden. Each also creates a new dialect, package culture, build system, error style, documentation habit, and set of failure modes that humans must learn before they can safely change a system.

The early languages compressed machine pain. Assembly gave names to opcodes and registers. Backus's Turing Award lecture framed the deeper trap as the ``von Neumann style'': languages and machines reinforcing each other's assignment-and-memory habits \cite{backusVonNeumann1978}. C made portable systems programming possible, as Ritchie's history of C makes clear, but it also made memory safety and low-level execution semantics permanent review burdens \cite{ritchieDevelopmentC1993}. Managed runtimes compressed memory management and deployment portability into VM policy. Scripting languages compressed iteration time and glue code, but often moved correctness from compiler feedback into convention, test discipline, and senior reviewer memory. The web multiplied the problem: JavaScript became the universal runtime by distribution, not by design, and TypeScript later compressed some of that chaos into a typed product surface \cite{typescript70Beta2026}. GitHub's 2025 Octoverse report describes TypeScript as the most-used language on GitHub and AI as mainstream in development \cite{githubOctoverse2025}; that is not an accident. The dominant language is now entangled with the dominant collaboration substrate.

The deepest advances were not only new syntax. Dijkstra's attack on unstructured control flow, Parnas's information-hiding decomposition criteria, and McCabe's complexity measure all tried to make local reasoning cheaper \cite{dijkstraGoto1968,parnasDecomposition1972,mccabeComplexity1976}. Conway's law warned that software structure also follows communication structure, so the codebase records organization as well as design \cite{conwayCommittees1968}. Those ideas still matter in agent-assisted work: a repository that hides ownership, boundary policy, and proof routes is asking the next author to reconstruct the design theory from residue.

There were also hopeful branches that did not become universal defaults. Julia attacked the two-language problem in scientific computing with multiple dispatch and JIT performance \cite{bezansonJulia2017}. Elm showed how frontend reliability can improve when impossible states are harder to express. F\#, Clojure, Haskell, Scala, Elixir, Nim, Crystal, D, and other ecosystems carried serious ideas about type systems, concurrency, supervision, native compilation, and expressiveness. Many survived as strong niches. Many failed to become default product infrastructure. The usual cause was not that the ideas were stupid. It was standard gravity: package maturity, deployment norms, hiring supply, build speed, debugging habits, cloud operations, security scanning, database migration practice, interop, and the amount of public example code that humans and agents can learn from.

This is the old human bottleneck. Developers do not merely learn syntax; they learn a language's social contract. Where does business truth live? Which errors are expected and which are exceptional? Which generated files are sacred? Which tests prove behavior and which only prove the mock? Which layer owns authorization? Which migration is safe? Which dependency is normal? Which code is legacy and which code is a pattern? Brooks's ``No Silver Bullet'' warned that the essential difficulty of software lives in the conceptual structure, not only the notation \cite{brooksNoSilverBullet1987}. In a large repository, that conceptual structure is rarely in one place. It is spread across conventions, old pull requests, tribal warnings, CI jobs, stale docs, and reviewers who remember why the last rewrite failed.

Technical debt names the same pressure over time. Cunningham's debt metaphor, Parnas's software-aging account, Lehman's evolution laws, and later technical-debt theory all describe how expedient decisions accumulate interest when the system keeps changing \cite{cunninghamTechnicalDebt1992,parnasSoftwareAging1994,lehmanEvolutionLaws1980,kruchtenTechnicalDebt2012,liTechnicalDebtMapping2015}. Refactoring is the repair promise, but empirical refactoring work shows that real teams refactor opportunistically and under imperfect knowledge, not as clean textbook phases \cite{brownManagingTechnicalDebt2010,murphyHillRefactoring2012}. In an AI-assisted repository, debt interest includes evidence interest: every hidden convention, missing owner route, vague test claim, and undocumented waiver raises the cost of the next safe change.

AI changes the economics without removing the mess. Languages reduced notation, runtime, deployment, and local reasoning burdens; AI compresses the next layer, the production of plausible edits. A model can imitate syntax, follow local style, and generate code in almost any mainstream language. That makes first drafts cheap. It does not make repository intent explicit, prove ownership, or bind a patch to merge evidence. It can even make hidden ambiguity more expensive, because the model confidently fills gaps that human teams used to negotiate slowly. DORA's 2025 research frames AI as an amplifier of existing organizational capability rather than an automatic productivity guarantee \cite{doraStateAIAssistedSoftwareDevelopment2025}. Security evidence is less forgiving: Veracode reports syntactically plausible generated code with known flaws \cite{veracodeGenaiCodeSecurity2025}, and GitGuardian reports accelerating secret exposure in public GitHub activity, including AI-assisted commits \cite{gitguardianSecretsSprawl2026,gitguardianSecretsSprawl2026Blog}.

The new bottleneck is not language expression. It is merge evidence. An agent opens a pull request that changes a checkout flow. The diff touches product UI, a generated API client, a schema-adjacent file, tests, and documentation. The hard reviewer question is no longer whether the patch looks plausible. It is: what evidence makes this safe to merge?

\begin{table}[t]
\caption{Bottleneck shift for agent-era repositories.}
\label{tab:bottleneck-shift}
\centering
\scriptsize
\begin{tabularx}{\columnwidth}{L{0.34\columnwidth} Y}
\toprule
\JKTableHeader
\textbf{Era} & \textbf{Primary bottleneck} \\
\midrule
Old bottleneck & Syntax, machine constraints, memory layout, and local comprehension. \\
Current bottleneck & Merge evidence, owner routing, proof freshness, generated-boundary integrity, and repair locality. \\
\bottomrule
\end{tabularx}
\end{table}

That question is the center of Jankurai. A vibe artifact is not merely bad code. It is repository residue that makes a change plausible but not auditable: an ownerless path, unmapped proof lane, hand-edited generated output, false-green test, missing security receipt, overbroad tool permission, unproven UI change, stale waiver, or review claim with no receipt. Vibe artifacts are the agent-era form of language chaos. They appear when the repository cannot say what owns the behavior, which proof applies, or what evidence would make the merge trustworthy.

The correct response is not to slow agents down until the old review process survives. The correct response is to redesign the repository as the alignment layer. The repository, not the prompt or reviewer memory, must encode ownership, allowed edit scope, proof lanes, generated boundaries, permission limits, evidence receipts, repair queues, and versioned conformance. Jankurai's operating claims are deliberately concrete: the repository is the alignment layer for human- and agent-authored code; vibe artifacts are detectable as repository states; findings are reducible when they route to owners, proof lanes, receipts, and repair queues; and agent efficiency improves when broad context rereads are replaced by path-scoped policy and reusable evidence.

\begin{figure*}[t]
\centering
\fbox{\begin{minipage}{0.94\textwidth}
\centering
\small
change proposal
\(\rightarrow\) changed paths
\(\rightarrow\) vibe artifact scan
\(\rightarrow\) owner route
\(\rightarrow\) proof lane
\(\rightarrow\) receipt
\(\rightarrow\) merge witness
\(\rightarrow\) \texttt{pass}/\texttt{review}/\texttt{block}/\texttt{ratchet\_fail}/\texttt{release\_fail}
\(\rightarrow\) repair queue or waiver
\(\rightarrow\) reduced baseline
\end{minipage}}
\caption{The Jankurai continuous repository alignment loop. A change is not trusted because it looks plausible; it is trusted only when changed paths, owners, proof receipts, missing evidence, artifact digests, and commit/tool identity are bound into a merge witness.}
\label{fig:alignment-loop}
\end{figure*}

The checkout-flow pull request is the running example. If the generated client changed by hand, the repository should emit \texttt{HLT-002-GENERATED-MUTATION}. If the UI changed without screenshot, ARIA, trace, or geometry evidence, it should emit \texttt{HLT-013-RENDERED-UX-GAP}. If the path has no owner or proof route, it should emit \texttt{HLT-003-OWNERLESS-PATH} or \texttt{HLT-004-UNMAPPED-PROOF}. The merge witness then records whether the PR can pass, needs human review, must block, violates a ratchet, or fails release gates.

Agent-native engineering is not ``humans approve whatever the model says.'' It is human-directed engineering in which intent becomes bounded authority, changed paths become proof lanes, proof lanes leave evidence, evidence drives repair or waiver expiry, and repeated repairs become reusable primitives. The same machinery must govern human-authored code, because a vague PR from a person and a vague patch from an agent fail in the same place: the repository cannot tell what owns the behavior, which proof applies, or what evidence would make the merge trustworthy.

This paper makes four contributions. First, it defines a stack-neutral repository conformance model for continuous repository alignment. Second, it specifies a merge-witness artifact that binds changed paths to ownership, proof, evidence, and decision. Third, it turns missing evidence into stable rule IDs and an ordered repair queue. Fourth, it reports current seed conformance-suite evidence and separates the non-normative reference profile from Jankurai Core.

The language story matters because it explains why another linter, another benchmark, or another prompt convention is not enough. Human teams already spent decades trying to make software safer by choosing better languages. That helped, but it never eliminated the repository-level questions. A Rust module can still hide the wrong owner. A TypeScript component can still ship a broken checkout state. A Python service can still own truth it should not own. A generated client can still be edited by hand. A CI suite can still be green without proving the changed behavior. The next professional baseline therefore cannot be only a language choice. It has to be a repository interface.

Jankurai is that interface. It does not ask every team to adopt the same stack or abandon existing tools. It asks every conformant repository to expose the evidence that serious review already needs: changed paths, owners, generated boundaries, required proof lanes, current receipts, missing evidence, artifact digests, waivers, repair actions, and a closed merge decision. The claim is simple but demanding: if a repository cannot explain why a change is safe to merge, the patch is still a vibe, no matter who or what wrote it.
